\section{Related Work \pete{WIP}}
\label{sec:related_work}

\cite{majumder2024data}
TBA.

\newpara{End-to-end Data-driven Discovery.} 
Most previous autonomous data-driven discovery systems, such as Bacon \citep{Langley1981DataDrivenDO, Langley1984TheSF, Langley1983RediscoveringCW} severely lacked the requisite computational power, restricting their scope with limited discovery of data-driven knowledge. A recent system, CoScientist \citet{Boiko2023AutonomousCR}, uses LGMs to automate some parts of the workflow; however, it still requires substantial human intervention (e.g., wet lab experiments) for hypothesis verification, thus not qualifying as a fully autonomous discovery system. DataLume \cite{Gu2023HowDD} fully automates the code generation for data transformation and hypothesis verification; however, do not support hypothesis search for complex science workflows. \citet{Gil2022TowardsCS, Gil2017TowardsCS, Gil2013UsingSW} as well as Automatic Analysis in WolframAlpha\footnote{\url{https://www.wolframalpha.com/examples/pro-features/data-input}} prototyped various workflows for conducting science in data-driven ways, however, such prototypes never explored the power of LGMs and are only exhibit limited generalizability to datasets and scientific methods. 

% CoScientist, Bacon, DataLume, Mathematica \harshit{Structure: Can be covered in this flow: [Bacon as the foundational system. Attempts by Mathematica. Then domain-specific attempts by Coscientist and DataLume.][Q: shall we also cover other attempts for doing science by DeepMind and MSFT at the start? I don't think any of them are end-to-end though. Or their coverage in the intro suffice? ]}

\newpara{AutoML.}
AutoML is a workflow of automatically building optimal machine learning and predictive models. AutoML tools exist in scientific packages like Scikit \cite{feurer2015efficient} and also in cloud platforms such as Google Cloud Platform\footnote{\url{cloud.google.com/automl}}, Microsoft Azure\footnote{\url{azure.microsoft.com/en-us/products/machine-learning/automatedml/}}, and Amazon Web Services\footnote{\url{aws.amazon.com/machine-learning/automl/}}. Existing AutoML Cloud platform systems primarily
% mainly focus on black box models, ensuring the models can be served at scale. Despite 
perform search over hyperparameters for optimal model development, but they cannot comprehend the semantics of the data and hence cannot help with data-driven hypothesis generation, orchestrating research pathways, and knowledge integration. MLAgentBench \cite{Huang2023BenchmarkingLL}, an evolution of AutoML, performs end-to-end machine learning to benchmark AI research agents. MLAgentBench can plan, evaluate hypotheses, and measure progress, but with a focus on optimizing machine learning models, not on discovering new and novel scientific knowledge.


% Google AutoML, MLAgentBench, AutoGluon/AWS Sagemaker, AutoML from other groups \harshit{Structure: AutoML tools exist in scientific packages like Scikit (link) and also in cloud platforms on Google Cloud Platform, Microsoft Azure and Amazon Web Services. Cloud platform systems like Google AutoML focus on black box models while ensuring the models can served at scale. MLAgentBench can be considered as an evolution of AutoML that does end-to-end ML to benchmark AI research agents. General AutoML tools can only help partially with hypothesis evaluation and measurement of progress. The cloud platforms tools can work on large datascale. But outside of this they do not help with hypothesis generation, planning and orchestrating research pathways and knowledge integration - something the ideal DDD system can do. MLAgentBench can plan, evaluate hypothesis and measure progress, but with a focus on optimal models, not for discovering new and novel science.}

\newpara{Automated Data Analysis.}
Automated Data Analysis tools are primarily focused on exploring data under a user-provided query (e.g., ``plot sales trends for last 12 months'', etc.) and often do not have the capability of searching through the hypotheses space as defined by the data.
% Tools such as PowerBI\footnote{\url{https://www.microsoft.com/en-us/power-platform/products/power-bi}}, Tableau\footnote{\url{https://www.tableau.com/}}, and Thoughtspot\footnote{\url{https://www.thoughtspot.com/}} can, however, perform multi-step hypothesis verification using inbuilt data transformation and statistical tools, though the interpretation and consumption of such analysis are left to the user.
Spreadsheet tools such as Microsoft Excel and Google Sheets are often part of the scientific workflow as the data analysis tool but show limited automation ability even after coding (Python) support \cite{PYExcel, duetai}. Focus on integrating LGMs into data analysis with known workflows \cite{Perlitz2022nBIIGAN, Chakraborty2024NavigatorAG}  and code-first data analysis \cite{ydata} increased recently---however, these are limited to small-scale tables and lack abilities such as orchestrating research plans, interpreting results, and knowledge integration. 